% Methodology Section for ICUFN 2026 Paper
% Context-Aware Drug Discovery with Zero-Fee Blockchain-Verified Biomaterial Provenance

\section{System Design and Methodology}

\subsection{System Architecture}

The proposed system integrates three core components: (1) BioPassport for biomaterial credential verification, (2) PureProtX for consensus AI-based drug screening, and (3) PureChain for zero-fee blockchain anchoring. Figure~\ref{fig:architecture} illustrates the overall architecture.

The system operates in three layers:
\begin{itemize}
    \item \textbf{Application Layer}: Handles user interactions, credential queries, and AI predictions
    \item \textbf{Verification Layer}: Performs policy checks, deterministic hashing, and provenance linking
    \item \textbf{Blockchain Layer}: Provides immutable storage via PureChain (Chain ID: 900520900520)
\end{itemize}

\subsection{Biomaterial Verification Protocol}

Before drug screening commences, the system verifies the biomaterial's provenance through BioPassport. A biomaterial is considered valid if and only if it satisfies the following predicate:

\begin{equation}
    \text{VALID} = \text{ID} \land \text{QC} \land \neg\text{Q} \land \neg\text{R} \land \text{TC}
\end{equation}

\noindent where ID = has IDENTITY credential, QC = has valid QC\_MYCO credential, Q = QUARANTINED status, R = REVOKED status, and TC = transfer chain is valid.

The verification process queries on-chain credential state and returns a credential hash $H_c$ computed as:

\begin{equation}
    H_c = \text{SHA256}(\text{canonical}(C))
\end{equation}

\noindent where $C$ represents the credential data and $\text{canonical}()$ produces deterministic JSON with sorted keys.

\subsection{Consensus AI Screening}

The AI screening module employs an ensemble of three machine learning models:

\begin{enumerate}
    \item Support Vector Regression (SVR) with RBF kernel
    \item Random Forest Regressor (100 estimators)
    \item Gradient Boosting Regressor (100 estimators)
\end{enumerate}

Each model receives a feature vector $\mathbf{x} \in \mathbb{R}^{2058}$ comprising:
\begin{itemize}
    \item 10 molecular descriptors (MW, LogP, HBD, HBA, TPSA, RotB, ArRings, Fsp3, HeavyAtoms, Charge)
    \item 2048-bit Morgan fingerprints (radius=2)
\end{itemize}

The consensus prediction is computed as:

\begin{equation}
    \text{pIC50}_{\text{consensus}} = \frac{1}{3}\sum_{m \in \{SVR, RF, GB\}} \text{pIC50}_m(\mathbf{x})
\end{equation}

\subsection{Unified Audit Record}

The core innovation is the Unified Audit Record (UAR), which links biomaterial provenance to screening results. The UAR schema contains:

\begin{itemize}
    \item Biomaterial provenance: material\_id, credential\_hash $H_c$, verification status
    \item Drug screening: molecule SMILES, model\_hash $H_m$, parameters\_hash $H_p$, results
    \item Master hash $H_{master}$ for single-point verification
\end{itemize}

The master hash is computed over the canonical JSON representation:

\begin{equation}
    H_{master} = \text{SHA256}(\text{canonical}(\{H_c, H_m, H_p, \text{results}, \text{metadata}\}))
\end{equation}

This ensures that any change in context (biomaterial, model, parameters, or results) produces a different hash.

\subsection{Blockchain Anchoring}

The master hash is anchored to PureChain, a zero-fee EVM-compatible blockchain using Proof of Authority consensus. Key properties:

\begin{itemize}
    \item \textbf{Zero transaction fees}: Enables unrestricted provenance logging
    \item \textbf{Deterministic finality}: PoA consensus provides fast confirmation
    \item \textbf{Independent verification}: Any party can verify by recomputing $H_{master}$
\end{itemize}

The anchoring transaction stores:
\begin{equation}
    \text{TX} = \{job\_id, molecule\_id, H_{master}, \text{metadata}\}
\end{equation}

\subsection{Context-Awareness Validation}

We define context-awareness through the following properties:

\begin{property}[Hash Uniqueness]
For any two distinct contexts $C_1 \neq C_2$:
\begin{equation}
    H_{master}(C_1) \neq H_{master}(C_2)
\end{equation}
\end{property}

\begin{property}[Deterministic Reproducibility]
For identical context $C$ across $N$ executions:
\begin{equation}
    |\{H_{master}(C)_1, H_{master}(C)_2, ..., H_{master}(C)_N\}| = 1
\end{equation}
\end{property}

These properties ensure that: (1) different experimental contexts produce distinguishable provenance records, and (2) identical experiments are perfectly reproducible.

\subsection{Experimental Protocol}

We evaluate the system through five experiment categories:

\begin{enumerate}
    \item \textbf{Context-Awareness Validation}: Verify that modifications to any context field produce distinct hashes
    \item \textbf{Deterministic Reproducibility}: Execute workflows $N$ times and verify 100\% hash match rate
    \item \textbf{Blockchain Performance}: Measure p50/p95/p99 latency and throughput
    \item \textbf{Provenance Completeness}: Document all captured artifacts and verification capabilities
    \item \textbf{Case Study}: Demonstrate end-to-end pipeline on sample drug molecules
\end{enumerate}

Performance metrics include:
\begin{itemize}
    \item Verification latency (target: p50 $<$ 50ms)
    \item Anchoring latency (target: p50 $<$ 100ms)
    \item Throughput (target: $>$ 10 tx/s)
    \item Gas cost: 0 (zero-fee blockchain)
\end{itemize}

\subsection{Implementation}

The system is implemented in Python 3.8+ using:
\begin{itemize}
    \item \textbf{Web3.py}: Blockchain interaction
    \item \textbf{scikit-learn}: Machine learning models
    \item \textbf{RDKit}: Molecular feature generation
    \item \textbf{hashlib}: SHA-256 hashing
\end{itemize}

The codebase comprises approximately 4,000 lines across core modules (PureProtX), integration modules (BioPassport), and experimental validation scripts.

% End of Methodology Section
